%% =====================================================================
%% §3 Method — 재작성 초안 v3 (2026-07-31)
%%
%% 대상: acmart-primary/mac_irl_icaif26.tex 280~746행을 이것으로 교체
%% 근거: docs/paper_ch3_rewrite_2026-07-31.md
%%
%% 반영 시 주의
%%  - Figure 1 (TikZ 파이프라인, 현행 292~417행) 이월 시 Reward & policy 박스를
%%    본문과 같은 순서로 바꿀 것: g 형태 → R_i → a*. 그리고 첨자
%%    $a^{*}_{i,t}$ → $a^{*}_{i,t+1}$.
%%  - §4.1 로 넘어가는 내용: CPCV 분할별 train/test 일수,
%%    Adam 설정(seed 42, 75 epochs, 유형별 batch·lr), 정확한 0 비율.
%% =====================================================================

\section{Method}
\label{sec:method}

%% [[FIGURE 1 — 현행 292~417행 TikZ 블록 그대로 이월]]

We compare investor types under a common specification. We first define
one-day-ahead net-buy behaviour, then the three candidate behavioural
channels and two market conditions that describe it, then a reward whose
optimal action generates that behaviour, and finally the same estimator
for every type. Four procedures and a feature ablation test whether the
recovered signs persist across resampling schemes, forecast orderings,
and estimators. We call this design shared-specification preference
recovery (SPR). Its coefficients describe conditional associations in
aggregate flow, not structural preferences.

\subsection{Problem setup}
\label{sec:method-setup}

Let $\mathcal{I} = \{\textrm{foreign}, \textrm{institutional},
\textrm{retail}\}$ index investor types. For type $i$ on day $t$, write
$V^{\mathrm{buy}}_{i,t}$ and $V^{\mathrm{sell}}_{i,t}$ for gross purchase
and sale value, and define the normalized net-buy action

\begin{equation}
  \label{eq:action}
  u_{i,t} = \frac{V^{\mathrm{buy}}_{i,t} - V^{\mathrm{sell}}_{i,t}}
                 {V^{\mathrm{buy}}_{i,t} + V^{\mathrm{sell}}_{i,t}}
          \in [-1, 1].
\end{equation}

Scaling net buying by the type's own gross traded value retains both
direction and intensity and places every type on the same scale
regardless of its trading size. The prediction target is
$a_{i,t+1} = u_{i,t+1}$, and the state $s_{i,t} = (x_{i,t}, C_t)$
contains only information observable through day~$t$. All three types
share the feature vector $x_{i,t} \in \mathbb{R}^3$, the market context
$C_t \in \mathbb{R}^2$, and the model form; only the parameters are
estimated separately, from each type's own flow. Each type is treated as
an agent acting independently, so no equilibrium restriction is imposed
and none is recovered. This symmetry makes the fitted coefficients
directly comparable across types without type-specific predictors.

\subsection{State features and market context}
\label{sec:method-features}

The state vector $x_{i,t} = (x^{\mathrm{mom}}_{t}, x^{\mathrm{per}}_{i,t},
x^{\mathrm{uw}}_{i,t})^{\top}$ represents three candidate channels:
recent price movement, the type's own flow persistence, and an aggregate
loss region. All three are specified in advance rather than selected by
search; each operationalizes one of the regularities reviewed in
Section~\ref{sec:related}.

\textbf{Recent price movement.} Let $P_t$ be the closing price. We use
the 20-day log return $x^{\mathrm{mom}}_t = \log(P_t / P_{t-20})$ with a
prespecified horizon~\cite{jegadeesh_titman}. A positive (negative)
coefficient is consistent with trend-following (contrarian)
behaviour~\cite{choe_kho_stulz,grinblatt_keloharju2001}; we expect the
former for foreign flow and the latter for retail flow, with no strong
prior for institutional flow.

\textbf{Own flow persistence.} $x^{\mathrm{per}}_{i,t} = a_{i,t-1}$, the
previous day's action on the same own-denominator definition
as~\eqref{eq:action}, so the coefficient is a first-order autoregression
on behaviour. Sias's decomposition of institutional
demand~\cite{sias2004} separates institutions following each other from
institutions following their own lagged trades; at the level of type
aggregates those two components cannot be separated, so we use only the
measurable own-following component and do not describe this feature as
herding. We expect a positive coefficient for all three types, following
the long-range dependence Oh~\cite{oh2025} reports for the netted daily
flow of each type in this market, and leave the cross-type ranking open,
since that paper finds the ordering does not survive netting.

\textbf{Aggregate loss region.} This feature approximates whether
retained inventory was acquired above the current trade price. Let $Q^b$
and $Q^s$ be gross purchase and sale quantity, $M^b$ and $M^s$ the
corresponding values, and $P^b = M^b / Q^b$ the purchase VWAP when
$Q^b > 0$; type and date subscripts are suppressed and
$[v]^{+} = \max(v, 0)$. With $\tilde H = \rho H_{t-1}$,

\begin{equation}
  \label{eq:underwater}
  \begin{aligned}
    O &= [\tilde H - Q^s]^{+}, &
    H &= [\tilde H + Q^b - Q^s]^{+}, &
    N &= H - O, \\
    \bar P &= \frac{O \bar P_{t-1} + N P^b}{H}, &
    P^{\mathrm{vwap}} &= \frac{M^b + M^s}{Q^b + Q^s},
  \end{aligned}
\end{equation}

where $\tilde H$ is decayed inventory entering day~$t$, $O$ the prior
inventory surviving sales, $N$ newly retained inventory, and
$H = O + N$ total effective inventory. The recursion updates both the
retained quantity and its average cost $\bar P$, while
$P^{\mathrm{vwap}}$ summarizes the same-day transaction price. We
initialize $H_{i,0} = 0$ and skip any VWAP term whose denominator is
zero. We define

\begin{equation}
  \label{eq:lossregion}
  x^{\mathrm{uw}}_{i,t} = \bigl[(\bar P_{i,t} - P^{\mathrm{vwap}}_{i,t})
    / \bar P_{i,t}\bigr]^{+},
\end{equation}

and set it to zero when effective inventory is zero or no trade occurs.
The forgetting factor $\rho = 0.98$ gives recent trades greater effective
weight, a half-life of roughly 34 trading days. This is an aggregate
proxy rather than an account-level cost basis; its sign can be compared
with disposition-effect predictions~\cite{shefrin_statman,odean1998} but
cannot identify that effect. We expect a positive coefficient for retail
flow, with no clear prior for the other two types.

\textbf{Market conditions.} Let $r_t$ be the one-day KOSPI~200 return and
$e_t$ the won--dollar closing rate. With $\mu_{t-1,252}$ and
$\sigma_{t-1,252}$ the mean and standard deviation of
$\{\log e_s : s = t-252, \dots, t-1\}$, define
$z^{\mathrm{FX}}_t = (\log e_t - \mu_{t-1,252}) / \sigma_{t-1,252}$ and
$C_t = (r_t, z^{\mathrm{FX}}_t)^{\top}$. Within each split, every feature
and context variable is standardized using training indices only, and the
resulting transformation is applied to the test data.

The three expectations above are fixed before estimation, so
Section~\ref{sec:experiments} reads the fitted coefficients against them
rather than rationalizing them afterwards.

\subsection{Reward and parameterization}
\label{sec:method-reward}

For type $i$ on day $t$ define the scalar

\begin{equation}
  \label{eq:gradient}
  g_{i,t} = (\beta_i + B_i C_t)^{\top} x_{i,t} + \alpha_i^{\top} C_t ,
\end{equation}

which separates an average feature sensitivity, a context-dependent
adjustment of that sensitivity, and a direct context effect. Here
$\beta_i \in \mathbb{R}^3$ is the sensitivity at mean context,
$B_i \in \mathbb{R}^{3 \times 2}$ how that sensitivity shifts with
context, and $\alpha_i \in \mathbb{R}^2$ the direct context effect when
the features are at their standardized means, so $\beta_i + B_i C_t$ is
the effective feature weight on day~$t$. Writing $\operatorname{rvec}$
for row-major vectorization and $[\,\cdot\,;\cdot\,]$ for vertical
concatenation,

\begin{equation}
  \label{eq:vec}
  \begin{aligned}
    w_{i,t} &= \bigl[\, x_{i,t};\ \operatorname{rvec}(x_{i,t} C_t^{\top});\ C_t \,\bigr], \\
    \theta_i &= \bigl[\, \beta_i;\ \operatorname{rvec}(B_i);\ \alpha_i \,\bigr]
      \in \mathbb{R}^{11},
  \end{aligned}
\end{equation}

so that $g_{i,t} = \theta_i^{\top} w_{i,t}$, with no intercept. We
interpret the five coefficients in $\beta_i$ and $\alpha_i$; the six
entries of $B_i$ are retained in every fit as conditioning terms.

Each type chooses its next-day action from the feasible set
$\mathcal{A} = [-1, 1]$, the range of the normalized net-buy
in~\eqref{eq:action}, under the reward

\begin{equation}
  \label{eq:reward}
  R_i(a \mid s_{i,t}) = g_{i,t}\, a - \tfrac{1}{2}\kappa a^2,
  \qquad \kappa > 0,
\end{equation}

so that $g_{i,t}$ is the marginal reward of the first unit traded and
$\kappa$ the curvature of the quadratic cost. $R_i$ is strictly concave
and $\mathcal{A}$ is compact, so the optimal action exists, is unique,
and is the projection of the unconstrained optimum onto $\mathcal{A}$:

\begin{equation}
  \label{eq:optimal}
  a^{*}_{i,t+1} = \arg\max_{a \in \mathcal{A}} R_i(a \mid s_{i,t})
  = \operatorname{clip}\!\bigl(g_{i,t}/\kappa,\, -1,\, 1\bigr),
\end{equation}

which we write $a^{*}_{i,t+1}(\theta)$ where the dependence on the
parameters matters. The clip is therefore the boundary of $\mathcal{A}$
rather than a link function chosen for convenience. Only the ratio
$g_{i,t}/\kappa$ is identified from observed actions, so we set
$\kappa = 1$ and read $g_{i,t}$ on the action scale; magnitudes are
comparable across types because every type's action lies on the common
$[-1,1]$ scale of~\eqref{eq:action}. Concavity is what makes intensity
informative: under a reward linear in $a$ the optimal action would sit at
a bound every day, and observed magnitudes would carry no information
about the state.

Because the horizon is one step, the recovered reward is informationally
equivalent to the conditional best response. We therefore read $\beta_i$
as a conditional association in aggregate flow rather than as a
structural preference, and do not separate preferences from beliefs or
institutional constraints: a type that buys after positive returns may
prefer trend exposure, may believe returns persist, or may operate under
a mandate that produces the same flow.

\subsection{Estimation and validation}
\label{sec:method-estimation}

The predicted action is~\eqref{eq:optimal} evaluated at $\hat\theta_i$,
and we fit $\theta_i$ by minimizing its squared error against the
observed action under an $L_1$ penalty. With $\mathcal{D}_i$ the training
indices for type~$i$,

\begin{equation}
  \label{eq:estimator}
  \hat\theta_i = \arg\min_{\theta}
  \Bigl\{ \tfrac{1}{|\mathcal{D}_i|}
    \sum_{t \in \mathcal{D}_i}
      \bigl( a^{*}_{i,t+1}(\theta) - a_{i,t+1} \bigr)^2
    + \lambda \lVert \theta \rVert_1 \Bigr\}.
\end{equation}

The saturation rate is zero for all three types in every split, so the
bound in~\eqref{eq:optimal} never binds and~\eqref{eq:estimator} reduces
to a Lasso in $\theta$. Across all 45 splits the design matrix
$W_i = [w_{i,t}^{\top}]_{t \in \mathcal{D}_i}$ has full column rank~11,
with a worst-case condition number of 6.94, so the reduced objective is
strictly convex and its minimizer is unique---which is what licenses
reading the reported coefficients as a property of the estimator rather
than of one solution among many.

The penalty is $\lambda = 0.005$ for every type. A common penalty is what
makes coefficient magnitudes comparable: under type-specific penalties a
smaller coefficient could reflect heavier shrinkage rather than weaker
sensitivity.

Four procedures and a feature ablation examine different sources of
instability, and Section~\ref{sec:experiments} attributes each claim to
the procedure that supports it.

\textbf{(V1) Combinatorial purged cross-validation.} We partition the
sample into ten chronological folds and take every choice of two folds as
the test set, giving $\binom{10}{2} = 45$
splits~\cite{lopezdeprado2018}, with a one-day purge at test boundaries
and a five-day embargo after each test block. This is the primary
protocol: it supplies the sampling distribution of $\beta_i$ and
$\alpha_i$ and the common basis for the ablation and penalty comparisons.
Because two test folds need not be adjacent, some splits train partly on
data later than their test block, so V1 measures sensitivity to
train--test composition rather than forecast skill, and we do not
interpret sign consistency as a $p$-value.

\textbf{(V2) Calendar-month block bootstrap.} Because the 45 splits share
observations, sign consistency across them understates uncertainty. We
resample calendar months with replacement 200 times, refitting each time,
and report 95\% intervals; monthly blocks preserve within-month
dependence in daily flow. This procedure fits in sample and produces no
held-out predictions, so we use it only for coefficient intervals.

\textbf{(V3) Expanding walk-forward evaluation.} Since V1 admits splits
that train on later data, we also impose strict chronology, training on
all prior calendar years and testing on 2023, 2024 and 2025 in turn, with
training sets of 242, 487 and 731 days. Scalers are fitted within each
window's training range.

\textbf{(V4) Penalty substitution.} To test dependence on the estimator
we refit under an $L_2$ penalty over a grid of fifteen strengths, select
one strength per type, and compare coefficient signs with the $L_1$ fits.

\textbf{Feature ablation.} We refit the model with each feature and each
context removed in turn, on the same 45 splits, and compare pooled
out-of-sample performance against the full specification using a paired
block bootstrap over dates with block length 20 and 10{,}000 resamples,
controlling the false discovery rate across variants and metrics. Pooling
across splits is required for the paired comparison, so these figures are
not directly comparable to the per-split averages reported for V1. We
report three single-feature and two context removals.

Performance is summarized by directional accuracy with
$\operatorname{sign}(0) = 0$, correlation between predicted and realized
actions, MAE, RMSE, and out-of-sample $R^2$ against each split's
training-mean action. Because action dispersion differs across types, we
do not rank types by absolute MAE.

Finally, we fix in advance what counts as identification. A coefficient
is reported as identified only when its sign agrees with its CPCV mean in
at least 90\% of the 45 V1 splits \emph{and} its V2 interval excludes
zero; otherwise it is reported as unidentified rather than interpreted.
V3 is reported alongside as evidence on temporal stability rather than on
identification: with only three windows the procedure is weak in both
directions. A feature whose removal does not degrade performance is
reported as contributing nothing.
